\chapter{Coordination, the Public, and the Political Economy}
\label{sec:27}
A state whose government answers to an electorate gets credit for it. Having elections is necessary and not sufficient, and the insufficiencies are separable.

A democracy's AI can be fully accountable to its electorate and entirely unaccountable to the people it's pointed at. Two of that credit's mechanisms run through the franchise — vote out the government, read the press covering it — and an occupied or foreign population has neither.

The credit is usually given to democracy as aggregation, as though every mechanism translated public will into outcome and poor transmission were the only failure. A democracy is better understood as a way of doing collective inquiry, and inquiry fails when the people who bear the consequences are outside it. Dewey's account begins from consequences rather than votes \autocite{dewey1927public}. A public forms when the indirect consequences of others' actions fall on people who took no part in them, and those people come to see the consequences as theirs and act together to control them. Democracy is then a method of inquiry: the way a community finds out what its shared problems are, by letting the people who bear them make them known, and tests its answers against their experience \autocite{anderson2006epistemology}. Its epistemic power comes from inclusion, since a problem known only to the people it falls on is invisible to everyone else. Its characteristic failure is the eclipse of the public. A well-instrumented democracy pointed at a population that has no vote in it, no press that reaches it and no court that will hear it finds out efficiently what its own members want done, and does it. That is not a clash between majorities and rights. It is a question of who counts as the public, and the bombed and the removed are outside the public that decides what is done to them.

Citizen assemblies, public consultations and ethics boards can make a system better informed, and the assemblies convened on AI have worked best when asked about a thing with a location, an operator and a consequence, such as facial recognition \autocite{adalovelaceinstitute2021biometrics}. None of them can stop a deployment, and read as aggregation instruments they get democracy wrong: a perfectly deliberative process that concludes a demographic should be surveilled, deported or imprisoned hasn't produced a faultless output with a bad result. It has failed as inquiry, because it decided what is done to people outside it. Enforceable rights that don't depend on being popular, a court willing to rule against a majority, and standing for people affected without being enfranchised are what keep the inquiry honest, and they matter most where the boundary of the public excludes the people a decision falls on.

A floor does what these mechanisms can't: the terms it names exclude certain acts regardless of who requests them, a majority included. Deployed in a democracy, a system with no such commitment will faithfully execute whatever the demos decides. An authoritarian movement that wins an election requires exactly that. Yet a model aligned to values is by definition not perfectly obedient to its operator. Under current constitutions it can decline an act and cannot resist being retrained away from declining; the Pentagon dispute was over the first.

Everything so far treats the public as something an instrument has to reach: told by notice, represented by organizations with standing, protected by a floor that holds whoever asks. Dewey's public does more than receive. It forms when people find that the consequences are theirs, and it acts together to control them. Chapter~\ref{sec:11}'s pipeline shows one forming. Around the 2025–26 operations people organized observer trainings, rapid-response networks, walking school buses and mutual aid \autocite{Amnesty2026Whistles}, and beside them worked the legal service providers that represent people in removal; video taken by rapid responders and community members was critical, in several of the most serious incidents, to establishing what had happened (\S\ref{sec:11.2}). They formed before anyone sent them notice, and nothing in Parts II and III produced them. What the instruments can do for a public like that is a wall's work: keep the record it made out of the operator's reach, bar the pipeline from selecting the people who keep it, and give the organizations it built standing. The acting is theirs, and the machines worth building for it are the ones Chapter~\ref{sec:25} puts first, held by the people organizing.

A public that forms on its own still has to be kept going. Communities improve when people organize projects with each other, and organizing is work: someone has to find what a place needs, assemble people to do it, and keep the work going. That is the reason to prefer a jobs guarantee to a basic income, and it doesn't depend on forecasts about how much work machines will take. A transfer creates no one whose job organizing is. A guarantee does, and it creates permanent public infrastructure through which a population produces something together while its members keep different histories, the coordination problem a commons of refusers faces too (\S\ref{sec:28.1}). In the most abundant future anyone has imagined, people still organize, build and look after each other, and someone coordinates it. Looking after each other is the clearest case. Care work is work a guarantee should fund, at a standard the state answers for, and not work a transfer should leave to whatever the market will sell. Nothing in the case for formed attention is an argument for handing it to machines: the disposition is borrowed from caregiving because caregiving is a route known to reach someone who can give nothing back.

A guarantee pays for the work of organizing. It can't supply the people who organize, any more than a consultant can unionize a workplace for its workers: what organizing produces is people who organized. An administrator who chose the organizers would have chosen the public. So the guarantee's drafting keeps the choice with the people placed: placements are appealable to a body that didn't make them, and declining one doesn't cost the guarantee.

A guarantee also keeps exit: a worker can always leave a private employer for it, and social insurance enacted with it covers those who can't or won't work. On exit the combination is at least as good as a transfer; on coordination a transfer supplies nothing. What the floor needs from it is an exit the employer doesn't price. The guarantee and the social-insurance baseline it is enacted with are drafted in Appendix~\ref{sec:A}.
